% from ega4-18-fr.tex
\oldpage[IV-4]{122}

$j:U\to X'$ d'un voisinage ouvert $U$ de $x$ dans $X$; restreignant au besoin $U$, on peut, en appliquant (I, 6.5.1, (i)), supposer que $h\circ j=f|U$, donc $j$ est de type fini (I, 6.3.4, (v)); enfin, l'homomorphisme $\varphi$ est surjectif par (18.4.6, (ii)), donc il résulte de (I, 6.5.4, (i)) qu'on peut, en restreignant encore $U$ et $X'$, supposer que $j$ est une immersion fermée. Cela prouve donc la nécessité de la condition énoncée; sa suffisance est immédiate (17.1.3, (i) et (ii)).

\begin{corollary}[18.4.8]
\label{IV.18.4.8}
Soit $f:X\to Y$ un morphisme localement de présentation finie. Pour que $f$ soit non ramifié (resp. étale), il faut et il suffit qu'il existe une famille de morphismes plats $g_\alpha:Y'_\alpha\to Y$ et, pour chaque $\alpha$, un ouvert $U_\alpha$ dans $X'_\alpha=X\times_Y Y'_\alpha$, tels que, si $g'_\alpha:X'_\alpha\to X$ et $f'_\alpha:X'_\alpha\to Y'_\alpha$ sont les projections canoniques, les $g'_\alpha(U_\alpha)$ forment un recouvrement de $X$, et que chacun des morphismes composés $U_\alpha\to X'_\alpha\xrightarrow{f'_\alpha}Y'_\alpha$ soit une immersion fermée (resp. ouverte). On peut de plus alors prendre les $g_\alpha$ étales.
\end{corollary}

La nécessité de la condition pour les morphismes étales est triviale, en prenant un seul $Y'_\alpha$ égal à $X$, le morphisme $g_\alpha$ correspondant étant égal à $f$, et l'ouvert $U\subset X\times_Y X$ étant la diagonale. Lorsque $f$ est non ramifié, la nécessité de la condition résulte de (18.4.7); on prend un recouvrement ouvert $(V_\alpha)$ de $X$ tel que pour chaque $\alpha$, $f|V_\alpha$ se factorise en $V_\alpha\xrightarrow{j_\alpha}Y'_\alpha\xrightarrow{g_\alpha}Y$, où $j_\alpha$ est une immersion fermée et $g_\alpha$ un morphisme étale. Alors $j_\alpha:V_\alpha\to Y'_\alpha$ se factorise en $V_\alpha\xrightarrow{s_\alpha}X'_\alpha\xrightarrow{f'_\alpha}Y'_\alpha$, où $s_\alpha$ est une $V_\alpha$-section de $X'_\alpha$, et comme le morphisme $g'_\alpha:X'_\alpha\to X$ est étale, $s_\alpha$ est une immersion ouverte (17.4.1), et il suffit de prendre $U_\alpha=s_\alpha(V_\alpha)$ pour répondre à la question.

La suffisance des conditions résulte de (17.7.1); on en conclut que $f$ est non ramifié (resp. étale) en chaque point de $g'(U_\alpha)$, donc dans $X$ tout entier puisque les $g'(U_\alpha)$ recouvrent $X$.

\begin{proposition}[18.4.9]
\label{IV.18.4.9}
Soient $S$ un préschéma, $f:X\to S$ un morphisme, $h:Y\to S$ un morphisme localement de présentation finie, $g:X\to Y$ un $S$-morphisme, $x$ un point de $X$, $y=g(x)$. Les conditions suivantes sont équivalentes:
\begin{enumerate}
  \item[a)] $h$ est étale au point $y$ et $g$ est plat au point $x$.
  \item[b)] $h$ est non ramifié au point $y$ et $f$ est plat au point $x$.
\end{enumerate}
\end{proposition}

Comme $f=h\circ g$, $a)$ entraîne que $f$ est plat au point $x$ (2.1.6), et évidemment que $h$ est non ramifié au point $y$, donc $a)$ implique $b)$.

Pour prouver que $b)$ entraîne $a)$, on peut d'abord supposer que $h$ est non ramifié (en remplaçant $Y$ par un voisinage de $y$); puis, en remplaçant $S$ par un voisinage ouvert de $s=h(y)=f(x)$, on peut supposer qu'il existe un morphisme étale $u:S'\to S$, un point $y'$ de $Y'=Y_{(S')}$ au-dessus de $y$ et un voisinage ouvert $V$ de $y'$ dans $Y'$ tel que, si $h'=h_{(S')}:Y'\to S'$, la restriction de $h'$ à $V$ soit une immersion fermée (18.4.8). Si l'on prouve alors que $h'$ est étale au point $y'$, il en résultera que $h$ est étale au point $y$ (17.7.1, (ii)); d'ailleurs, $f'=f_{(S')}:X_{(S')}\to S'$ est plat en tout point $x'$ au-dessus de $x$. Comme les projections $v:Y'\to Y$, $w:X'\to X$ sont des morphismes étales (donc plats), si l'on prouve que $g'=g_{(S')}:X_{(S')}\to Y_{(S')}$ est plat au point $x'$, il résultera de (2.2.11, (iv)) que $g$ sera plat au point $x$. On peut donc se borner au cas où $h$ est une \emph{immersion fermée}

